% =============================================================
%  Module Description — Calculus I (M102-S27)
%  Shown on https://wosm.academy as the course page for M102-S27.
%  Compile with: pdflatex m102-module-description.tex
%  Write "TODO(owner): ..." where information is still missing;
%  the website highlights it.
% =============================================================
\documentclass[11pt,a4paper]{article}

\usepackage[utf8]{inputenc}
\usepackage[T1]{fontenc}
\usepackage{lmodern}
\usepackage[margin=2.2cm]{geometry}
\usepackage{amssymb}
\usepackage{booktabs}
\usepackage{tabularx}
\usepackage{array}
\usepackage{enumitem}
\usepackage{xcolor}
\usepackage{titlesec}
\usepackage{fancyhdr}
\usepackage[hidelinks]{hyperref}

% ---------- Module details (edit these) ----------------------
% \ModuleCode must equal the course's `code` in src/content/courses/.
\newcommand{\ModuleTitle}{Calculus I}
\newcommand{\ModuleCode}{M102-S27}
\newcommand{\Programme}{Foundation programme}
\newcommand{\Level}{Semester 2}
\newcommand{\Semester}{Spring 2027}
\newcommand{\Credits}{7 ECTS}
\newcommand{\ModuleLeader}{Dr. Ali Khudiyev}
\newcommand{\Email}{ali.khudiyev@wosm.academy}
\newcommand{\OfficeHours}{Monday--Friday, 10:00--17:00 (Baku time, UTC+4)}
\newcommand{\Prerequisites}{Discrete Structures (M101-F26)}

% ---------- Styling ------------------------------------------
\definecolor{accent}{RGB}{0,70,130}
\titleformat{\section}{\large\bfseries\color{accent}}{\thesection.}{0.5em}{}
\titlespacing*{\section}{0pt}{1.2em}{0.5em}
\setlist{itemsep=2pt, topsep=4pt}
\setlength{\parindent}{0pt}
\setlength{\parskip}{4pt}

\pagestyle{fancy}
\fancyhf{}
\lhead{\small \ModuleCode\ -- \ModuleTitle}
\rhead{\small Module Description}
\cfoot{\small \thepage}
\renewcommand{\headrulewidth}{0.4pt}

% =============================================================
\begin{document}

\begin{center}
  {\LARGE\bfseries\color{accent} \ModuleTitle}\\[4pt]
  {\large Module Description}
\end{center}
\vspace{6pt}

% ---------- Summary table (not shown on the website: the course page
% shows these details from the macros above) --------------------------
\begin{tabularx}{\textwidth}{>{\bfseries}l X >{\bfseries}l X}
\toprule
Module code   & \ModuleCode    & Credits      & \Credits \\
Programme     & \Programme     & Level        & \Level \\
Semester      & \Semester      & Prerequisites & \Prerequisites \\
Module leader & \ModuleLeader  & Email        & \href{mailto:\Email}{\Email} \\
Office hours  & \multicolumn{3}{l}{\OfficeHours} \\
\bottomrule
\end{tabularx}

% ---------- Sections (shown on the website) ------------------
\section{Module Overview}
Calculus I is the first module of analysis in the curriculum. It builds the
calculus of functions of one real variable on a precise footing: the real numbers,
limits and continuity come first, and differentiation and integration are then
developed from them rather than taken as a set of rules.

The module covers:
\begin{itemize}
  \item The real numbers, sequences and limits (first 4--5 weeks)
  \item Continuity, differentiation and its applications (optimisation, approximation,
        Newton's method)
  \item The Riemann integral, the Fundamental Theorem of Calculus and techniques of
        integration
\end{itemize}

Throughout, the module links the theory to computer science: convergence of
iterative algorithms, asymptotic growth of functions, numerical approximation and
floating-point error. It prepares students for Calculus II and for Probability
Theory and Statistics.

\section{Aims}
\begin{itemize}
  \item Give a rigorous understanding of limits, continuity, derivatives and integrals of functions of one real variable.
  \item Develop fluency in computing derivatives and integrals and in using them to model and solve problems.
  \item Train students to read and write short, correct proofs in analysis using $\varepsilon$--$\delta$ arguments.
  \item Show how calculus underpins algorithms and numerical methods used in computer science.
\end{itemize}

\section{Learning Outcomes}
On successful completion of this module, students will be able to:
\begin{enumerate}[label=LO\arabic*.]
  \item State the completeness property of the real numbers and use it to reason about suprema, infima and convergence of sequences.
  \item Compute limits of sequences and functions, and prove simple limit and continuity statements with $\varepsilon$--$\delta$ arguments.
  \item Differentiate elementary functions and apply the Mean Value Theorem, L'H\^opital's rule and Taylor polynomials to analyse and approximate functions.
  \item Solve optimisation and related-rates problems and use Newton's method and bisection to approximate roots, estimating the error.
  \item Evaluate definite and indefinite integrals using the Fundamental Theorem of Calculus, substitution, integration by parts and partial fractions.
  \item Use integrals to compute areas, volumes and averages, and compare the growth of functions as used in asymptotic analysis of algorithms.
\end{enumerate}

\section{Syllabus}
\begin{tabularx}{\textwidth}{c X}
\toprule
\textbf{Unit} & \textbf{Topic} \\
\midrule
1  & The real numbers: ordered fields, absolute value, inequalities \\
2  & Completeness: supremum, infimum, the Archimedean property \\
3  & Sequences: convergence, limit laws, bounded and monotone sequences \\
4  & Subsequences, the Bolzano--Weierstrass theorem, Cauchy sequences \\
5  & Functions: elementary functions, inverses, exponential and logarithm \\
6  & Limits of functions: $\varepsilon$--$\delta$ definition, one-sided limits \\
7  & Continuity and its algebra \\
8  & Intermediate Value Theorem, Extreme Value Theorem; the bisection method \\
9  & Uniform continuity \\
10 & The derivative: definition, rules, chain rule \\
11 & Derivatives of inverse, exponential, logarithmic and trigonometric functions \\
12 & Mean Value Theorem, Rolle's theorem, monotonicity \\
13 & L'H\^opital's rule; growth rates and comparison with $O$-notation \\
14 & Higher derivatives, convexity, curve sketching \\
15 & Optimisation and related rates \\
16 & Newton's method and its convergence \\
17 & Taylor polynomials and Taylor's theorem with remainder \\
18 & The Riemann integral: partitions, upper and lower sums \\
19 & Integrability of continuous and monotone functions; properties of the integral \\
20 & The Fundamental Theorem of Calculus \\
21 & Integration by substitution and by parts \\
22 & Partial fractions and trigonometric integrals \\
23 & Improper integrals \\
24 & Applications: area, volume, arc length, average value \\
25 & Numerical integration: trapezoidal and Simpson's rules, error bounds \\
26 & Sums and integrals: bounding $\sum f(k)$ by $\int f$, harmonic numbers \\
\bottomrule
\end{tabularx}

\section{Teaching and Learning Methods}
The module is self-paced. All materials are provided on Moodle: lecture notes and
recorded lectures for each unit, tutorial problem sheets with worked solutions,
short quizzes, and the briefs for the two projects. Problem sheets are for practice
and are not marked. Questions are discussed in the module channel on Discord and in
office hours. Independent study includes reading, problem sheets and project work.

\begin{tabularx}{0.6\textwidth}{X r}
\toprule
\textbf{Activity} & \textbf{Hours} \\
\midrule
Lectures            & 52 \\
Labs / Tutorials    & 28 \\
Independent study   & 130 \\
\midrule
\textbf{Total}      & \textbf{210} \\
\bottomrule
\end{tabularx}

\section{Assessment}
Knowledge is assessed with weekly unit quizzes, two projects and a final
exam. Quizzes check computational fluency. The projects combine short proofs,
calculation and a small amount of code:
\begin{itemize}
  \item \textbf{Project 1 -- Roots and approximation:} implement bisection and
        Newton's method, compare their convergence, and use Taylor polynomials to
        approximate functions with proven error bounds.
  \item \textbf{Project 2 -- Integrals and asymptotics:} implement the trapezoidal
        and Simpson's rules with error bounds, and use integrals to bound sums such as
        harmonic numbers that appear in the analysis of algorithms.
\end{itemize}
The final exam tests whether students can choose and apply the right technique and
justify each step.

\begin{tabularx}{\textwidth}{X c c c}
\toprule
\textbf{Component} & \textbf{Weight} & \textbf{Due} & \textbf{LOs} \\
\midrule
Unit quizzes (best 10 of 12) & 20\% & Weekly & LO1--LO3, LO5 \\
Project 1: Roots and approximation & 20\% & Week 7 & LO2--LO4 \\
Project 2: Integrals and asymptotics & 20\% & Week 13 & LO5, LO6 \\
Final exam (3 hours, written) & 40\% & End of term & LO1--LO6 \\
\midrule
\textbf{Total}      & \textbf{100\%} & & \\
\bottomrule
\end{tabularx}

\section{Reading List}
\textbf{Essential}
\begin{itemize}
  \item M.\ Spivak: \textit{Calculus}, 4th ed., Publish or Perish, 2008.
  \item G.\ Strang, E.\ Herman: \textit{Calculus, Volume 1}, OpenStax, 2016. Free online at \url{https://openstax.org/details/books/calculus-volume-1}.
\end{itemize}
\textbf{Recommended}
\begin{itemize}
  \item J.\ Stewart, D.\ Clegg, S.\ Watson: \textit{Calculus: Early Transcendentals}, 9th ed., Cengage, 2020.
  \item T.\ M.\ Apostol: \textit{Calculus, Volume 1}, 2nd ed., Wiley, 1967.
  \item S.\ Abbott: \textit{Understanding Analysis}, 2nd ed., Springer, 2015.
  \item R.\ L.\ Graham, D.\ E.\ Knuth, O.\ Patashnik: \textit{Concrete Mathematics}, 2nd ed., Addison-Wesley, 1994 (Chapter 9, asymptotics).
\end{itemize}

\section{Materials and Licence}
Course materials are published under the
\href{https://creativecommons.org/licenses/by-sa/4.0/}{Creative Commons Attribution-ShareAlike 4.0}
licence unless stated otherwise. See \href{https://wosm.academy/licence/}{wosm.academy/licence} for details.

\section{Policies}
\begin{itemize}
  \item \textbf{Attendance:} Not mandatory. The module is self-paced; work through the units at your own pace within the term.
  \item \textbf{Late submissions:} Accepted, but penalised.
  \item \textbf{Academic integrity:} All submitted work must be your own.
        See the \href{https://wosm.academy/code-of-conduct/}{Code of Conduct}.
\end{itemize}

\end{document}
